\begin{proof}[Proof of \cref{obj:thm:heavy-tail-comparison}]
\begin{enumerate}
\item Fix an admissible public tuple and its matched smoothness tuple:
\[
v=(p,\alpha,\beta,\gamma)\in\mathcal V,
\qquad
w=(\alpha,\beta,\gamma),
\qquad
v_2=(2,\alpha,\beta,\gamma)\in\mathcal V.
\]
Use the quantities of \cref{obj:def:sharp-frontier-scales}, in particular
\[
q=\frac{p-1}{p},
\qquad
S=\alpha+\beta,
\qquad
E(p)=\min\{E_0(p),E_4(p)\}.
\]
Admissibility gives
\[
p-1>0,\qquad
0<q\le\frac12,\qquad
S>0,\qquad
\gamma>0,\qquad
2\gamma+q>0,\qquad
D_p>0.
\]
Indeed, \(q\le1/2\) follows by multiplying by \(p>0\) and using \(p\le2\); every term added to \(1\) in \(D_p\) is positive. Consequently,
\[
E_0(p)>0,\qquad E_4(p)>0,\qquad E(p)>0.
\]
The same conclusions apply at \(v_2\).

We first establish the phase selection needed for the lower bounds. Since
\[
pq=p-1,
\qquad
\frac1q-1=\frac1{p-1},
\]
subtracting the two channel exponents gives
\[
\begin{aligned}
E_4(p)-E_0(p)
&=\frac{2S(2\gamma+q)-2\gamma qD_p}
        {(2\gamma+q)D_p}\\
&=\frac{2\gamma q}{(2\gamma+q)D_p}
  \left(
  \frac{2S}{q}-2\alpha-\frac{\beta}{q}
  +\frac{S}{2\gamma}-1
  \right)\\
&=\frac{2\gamma q}{(2\gamma+q)D_p}\bigl(F(p)-1\bigr).
\end{aligned}
\]
The coefficient multiplying \(F(p)-1\) is strictly positive. Thus
\[
\begin{aligned}
F(p)\ge1&\ \Longrightarrow\ E_0(p)\le E_4(p)
                       \ \Longrightarrow\ E(p)=E_0(p),\\
F(p)<1&\ \Longrightarrow\ E_4(p)<E_0(p)
                       \ \Longrightarrow\ E(p)=E_4(p).
\end{aligned}
\]
These implications also hold with \(p=2\).

The definition of the lower multiplier gives
\[
N_{\mathrm{low}}(v)\ge2,
\qquad
c_0^{\mathrm e}>0,
\qquad
k_v>0,
\qquad
c_v>0,
\]
because every power occurring in either branch has a positive base. Applying this argument at \(v_2\), and applying \cref{obj:thm:original-record-attainable-rate} at both tuples, yields
\[
c_w^{\mathrm b}>0,\qquad C_v>0,\qquad C_w^{\mathrm b}>0.
\]
For every \(n\ge2\), positivity of real powers also gives
\[
\rho_n(v)>0,
\qquad
\rho_n^{\mathrm b}(w)>0.
\]
% lean: CausalSmith.Stat.FinitepHomogeneityDensegamma.phase_channel_difference
% lean: CausalSmith.Stat.FinitepHomogeneityDensegamma.phase_channel_selection
% lean: CausalSmith.Stat.FinitepHomogeneityDensegamma.cLower_pos
% lean: CausalSmith.Stat.FinitepHomogeneityDensegamma.matched_scales_pos

\item Fix an integer \(n\ge2\). We obtain the bounded lower bound by splitting according to the phase at \(v_2\).

If \(F(2)<1\), then
\[
E(2)=E_4(2),
\qquad
c_w^{\mathrm b}\rho_n^{\mathrm b}(w)
=c_w^{\mathrm b}n^{-E_4(2)}.
\]
The bounded conclusion of \cref{obj:lem:rough-full-record-sharp-lower} supplies normalized finite alternative witnesses and
\[
B_n^{\mathrm b}
 \bigl(w,c_w^{\mathrm b}\rho_n^{\mathrm b}(w)\bigr)
\ge\frac25.
\]
A normalized finite prior has a positive-weight atom: otherwise every nonnegative weight would be zero, contradicting that their sum is one. Choosing such an alternative atom gives
\[
\exists P\in\mathcal M_w^{\mathrm b}:
\quad
c_w^{\mathrm b}\rho_n^{\mathrm b}(w)\le d(P).
\]

If \(F(2)\ge1\), the public definitions and the phase selection give
\[
E(2)=E_0(2)=\frac{2\gamma}{4\gamma+1},
\qquad
c_w^{\mathrm b}=c_0^{\mathrm e}
=\frac{4^{-\gamma}}{16\sqrt3}.
\]
The signed-binary witnesses in \cref{obj:lem:paired-tent-full-record-lower} therefore give the same risk bound and a positive-weight alternative satisfying
\[
P\in\mathcal M_w^{\mathrm{bin}}
\subseteq\mathcal M_w^{\mathrm b},
\qquad
c_w^{\mathrm b}\rho_n^{\mathrm b}(w)\le d(P),
\]
where the inclusion follows from \cref{obj:prop:bounded-submodel-comparison}. Hence, in both cases,
\[
\exists P\in\mathcal M_w^{\mathrm b}:
\quad c_w^{\mathrm b}\rho_n^{\mathrm b}(w)\le d(P),
\qquad
B_n^{\mathrm b}
 \bigl(w,c_w^{\mathrm b}\rho_n^{\mathrm b}(w)\bigr)
\ge\frac25.
\]

We also verify that this bounded lower separation lies below the bounded cap. For any admissible tuple, the branch \(F(p)\ge1\) gives \(c_v=c_0^{\mathrm e}\). In the other branch, the phase identity gives \(E_0(p)-E_4(p)>0\), so
\[
N_{\mathrm{low}}(v)^{-(E_0(p)-E_4(p))}\le1,
\qquad
c_v\le
c_0^{\mathrm e}N_{\mathrm{low}}(v)^{-(E_0(p)-E_4(p))}
\le c_0^{\mathrm e}.
\]
Since \(E(p)>0\) and \(n\ge2\),
\[
0<\rho_n(v)\le1,
\qquad
0<c_v\rho_n(v)\le c_0^{\mathrm e}
\le\frac1{16\sqrt3}
<\frac1{8\sqrt2}=d_0.
\]
The strict numerical inequality follows from
\(\sqrt2<2\sqrt3\), obtained by squaring these positive quantities. Applying this bound at \(v_2\), and using
\(d_0\le D_w^{\mathrm b}\) from \cref{obj:prop:bounded-submodel-comparison}, gives
\[
0<c_w^{\mathrm b}\rho_n^{\mathrm b}(w)
<d_0\le D_w^{\mathrm b}.
\]

For the original model, \cref{obj:thm:tail-essential-full-record-converse} supplies
\[
0<c_v\rho_n(v)<D_v,
\qquad
B_n\bigl(v,c_v\rho_n(v)\bigr)\ge\frac25,
\]
and a normalized finite alternative prior. Choosing a positive-weight atom as above gives
\[
\exists P\in\mathcal M_v:
\quad c_v\rho_n(v)\le d(P)<D_v.
\]

We now pass from these failures to the capped radii. For
\(0<r_1\le r_2<D_v\), with a nonempty alternative class at \(r_2\), inclusion of the alternative classes gives, for every admissible test,
\[
\sup_{\substack{P\in\mathcal M_v\\ d(P)\ge r_2}}
 \{1-\mathsf R_n(P,\phi)\}
\le
\sup_{\substack{P\in\mathcal M_v\\ d(P)\ge r_1}}
 \{1-\mathsf R_n(P,\phi)\}.
\]
The errors lie in \([0,1]\), and the null-size constraint is unchanged. Taking infima over tests therefore yields
\[
B_n(v,r_2)\le B_n(v,r_1).
\]
The identical argument gives bounded-risk monotonicity.

If a successful separation \(\widetilde r\in(0,D_v)\) satisfied
\(\widetilde r<c_v\rho_n(v)\), monotonicity and its success condition would imply
\[
\frac25
\le B_n\bigl(v,c_v\rho_n(v)\bigr)
\le B_n(v,\widetilde r)
\le\frac1{10},
\]
a contradiction. The sentinel \(D_v\) also exceeds \(c_v\rho_n(v)\). Thus every element of the nonempty set defining the capped infimum in \cref{obj:def:critical-radius} is at least \(c_v\rho_n(v)\). The bounded argument uses its positive-weight witness and the legal separation established above. Consequently,
\[
c_v\rho_n(v)\le r_n^*(v),
\qquad
c_w^{\mathrm b}\rho_n^{\mathrm b}(w)
\le r_n^{*,\mathrm b}(w).
\]
The displayed risk bounds also imply both asserted strict inequalities against \(1/10\).
% lean: CausalSmith.Stat.FinitepHomogeneityDensegamma.bounded_matched_lower
% lean: CausalSmith.Stat.FinitepHomogeneityDensegamma.matched_lower_separation
% lean: CausalSmith.Stat.FinitepHomogeneityDensegamma.normalized_prior_positive_atom
% lean: CausalSmith.Stat.FinitepHomogeneityDensegamma.cappedRadius_lower_of_nonempty_failure

\item By \cref{obj:thm:original-record-attainable-rate}, the public tests satisfy
\[
r_n^*(v)\le C_v\rho_n(v),
\qquad
r_n^{*,\mathrm b}(w)\le C_w^{\mathrm b}\rho_n^{\mathrm b}(w).
\]
Together with \cref{obj:prop:bounded-submodel-comparison} and the preceding lower bounds, this gives
\[
\begin{aligned}
c_v\rho_n(v)
&\le r_n^*(v)\le C_v\rho_n(v),\\
c_w^{\mathrm b}\rho_n^{\mathrm b}(w)
&\le r_n^{*,\mathrm b}(w)
=r_n^{*,\mathrm{bin}}(w)
\le C_w^{\mathrm b}\rho_n^{\mathrm b}(w).
\end{aligned}
\]
The same attainment result supplies
\[
\begin{aligned}
\mathsf R_n(P,\phi^*_{n,v})&\le0.0075<\frac1{10}
&&\quad(P\in H_0(v)),\\
\mathsf R_n(P,\phi^{*,\mathrm b}_{n,w})&\le0.0075<\frac1{10}
&&\quad(P\in H_0^{\mathrm b}(w)).
\end{aligned}
\]
Under the respective conditions
\[
C_v\rho_n(v)<D_v,
\qquad
C_w^{\mathrm b}\rho_n^{\mathrm b}(w)<D_w^{\mathrm b},
\]
it supplies, respectively,
\[
\begin{aligned}
1-\mathsf R_n(P,\phi^*_{n,v})&\le0.0075<\frac1{10}
&&\quad(P\in\mathcal M_v,\ d(P)\ge C_v\rho_n(v)),\\
1-\mathsf R_n(P,\phi^{*,\mathrm b}_{n,w})&\le0.0075<\frac1{10}
&&\quad(P\in\mathcal M_w^{\mathrm b},\
d(P)\ge C_w^{\mathrm b}\rho_n^{\mathrm b}(w)).
\end{aligned}
\]

The bounded radius is positive by its lower bound. Dividing the matched bounds therefore gives
\[
\frac{c_v\rho_n(v)}
     {C_w^{\mathrm b}\rho_n^{\mathrm b}(w)}
\le\Delta_n(v)
\le
\frac{C_v\rho_n(v)}
     {c_w^{\mathrm b}\rho_n^{\mathrm b}(w)}.
\]
For \(n>0\), the law of real powers gives
\[
\frac{\rho_n(v)}{\rho_n^{\mathrm b}(w)}
=\frac{n^{-E(p)}}{n^{-E(2)}}
=n^{E(2)-E(p)}.
\]
Hence
\[
\frac{c_v}{C_w^{\mathrm b}}n^{E(2)-E(p)}
\le\Delta_n(v)
\le
\frac{C_v}{c_w^{\mathrm b}}n^{E(2)-E(p)}.
\]
Finally, \cref{obj:prop:bounded-submodel-comparison} gives
\[
0<r_n^{*,\mathrm b}(w)\le r_n^*(v),
\qquad
1\le\Delta_n(v).
\]
% lean: CausalSmith.Stat.FinitepHomogeneityDensegamma.matched_radius_ratio_bounds
% lean: CausalSmith.Stat.FinitepHomogeneityDensegamma.frontier_at

\item Since \(E(p)>0\) and \(E(2)>0\), the definitions of the scales give
\[
\rho_n(v)=n^{-E(p)}\longrightarrow0,
\qquad
\rho_n^{\mathrm b}(w)=n^{-E(2)}\longrightarrow0.
\]
The public thresholds are those of \cref{obj:def:sharp-frontier-scales}, evaluated at \(v\) and \(v_2\). Their formulas in \cref{obj:thm:original-record-attainable-rate} give
\[
N_v\ge4,\qquad N_w^{\mathrm b}\ge4.
\]
Its witness-threshold conclusion therefore applies to every natural number satisfying the corresponding threshold and yields
\[
\begin{aligned}
n\ge N_v&\ \Longrightarrow\ C_v\rho_n(v)<d_0,\\
n\ge N_w^{\mathrm b}&\ \Longrightarrow\
C_w^{\mathrm b}\rho_n^{\mathrm b}(w)<d_0.
\end{aligned}
\]
% lean: CausalSmith.Stat.FinitepHomogeneityDensegamma.rho_tendsto_zero
% lean: CausalSmith.Stat.FinitepHomogeneityDensegamma.rhoBounded_tendsto_zero
% lean: CausalSmith.Stat.FinitepHomogeneityDensegamma.frontier_at

\item We make the exponent difference explicit, always preserving \(w\) when evaluating at \(2\). Define
\[
D_2=1+2S+\frac{S}{2\gamma},
\qquad
F_2=2S+\frac{S}{2\gamma}.
\]
Substitution of \(q=(p-1)/p\) gives
\[
D_p=D_2+\frac{\beta(2-p)}{p-1},
\qquad
F(p)=F_2+\frac{(2\alpha+\beta)(2-p)}{p-1}
\ge F_2.
\]
For the supplied-propensity channel, a common denominator yields
\[
\begin{aligned}
E_0(2)-E_0(p)
&=\frac{2\gamma}{4\gamma+1}
 -\frac{2\gamma(p-1)}{2\gamma p+p-1}\\
&=\frac{4\gamma^2(2-p)}
 {(4\gamma+1)(2\gamma p+p-1)}.
\end{aligned}
\]
Both denominators are positive. For the rough channel, the identity for \(D_p\) gives
\[
\begin{aligned}
E_4(2)-E_4(p)
&=\frac{2S}{D_2}-\frac{2S}{D_p}\\
&=\frac{2S(D_p-D_2)}{D_2D_p}\\
&=\frac{2S\beta(2-p)}{(p-1)D_pD_2}.
\end{aligned}
\]
Combining these identities with the phase rules \(E(p)=E_0(p)\) for \(F(p)\ge1\) and \(E(p)=E_4(p)\) for \(F(p)<1\), applied also at \(2\), gives the exhaustive expression
\[
E(2)-E(p)=
\begin{cases}
\displaystyle
\frac{4\gamma^2(2-p)}
 {(4\gamma+1)(2\gamma p+p-1)},
&F_2\ge1,\\[6pt]
\displaystyle
E_4(2)-E_0(p),
&F_2<1\le F(p),\\[4pt]
\displaystyle
\frac{2S\beta(2-p)}{(p-1)D_pD_2},
&F(p)<1.
\end{cases}
\]
Indeed, in the first case \(F(p)\ge F_2\ge1\), so both minimum exponents select the supplied-propensity channel. In the second case the bounded exponent selects the rough channel and the finite-moment exponent selects the supplied-propensity channel. In the third case \(F_2\le F(p)<1\), so both select the rough channel.

If \(p<2\), the two channel differences displayed above are strictly positive. Therefore
\[
E(p)\le E_0(p)<E_0(2),
\qquad
E(p)\le E_4(p)<E_4(2),
\]
and hence
\[
E(p)<\min\{E_0(2),E_4(2)\}=E(2).
\]
Using the scale definitions in \cref{obj:def:sharp-frontier-scales}, we obtain
\[
\frac{\rho_n(v)}{\rho_n^{\mathrm b}(w)}
=n^{E(2)-E(p)}
\longrightarrow\infty.
\]
% lean: CausalSmith.Stat.FinitepHomogeneityDensegamma.phase_bounded_identities
% lean: CausalSmith.Stat.FinitepHomogeneityDensegamma.phase_bounded_tail_difference
% lean: CausalSmith.Stat.FinitepHomogeneityDensegamma.phase_bounded_rough_difference
% lean: CausalSmith.Stat.FinitepHomogeneityDensegamma.phase_difference_eq
% lean: CausalSmith.Stat.FinitepHomogeneityDensegamma.exponent_phase_algebra

\item Suppose \(p<2\). Apply \cref{obj:thm:tail-essential-full-record-converse} to the selected priors and magnitudes of \cref{obj:def:converse-handle}. It gives
\[
L_{n,v}\longrightarrow\infty.
\]
For every \(n\ge2\), both selected priors are normalized finite probability priors, and their positive-weight laws satisfy
\[
\begin{aligned}
\pi_{0,n,v}(\{P\})>0
&\ \Longrightarrow\ P\in H_0(v),\\
\pi_{1,n,v}(\{P\})>0
&\ \Longrightarrow\
P\in\mathcal M_v,\quad
c_v\rho_n(v)\le d(P)<D_v.
\end{aligned}
\]
For each positive-weight law in either prior, their conditional outcome kernels satisfy
\[
Q_{a,P}\bigl(\{0,-L_{n,v},L_{n,v}\}\mid x\bigr)=1
\]
for \(a\in\{0,1\}\) and \(\lambda\)-almost every \(x\). The same result supplies the complete-record bound
\[
\operatorname{TV}(\mathbb P_{0,n,v},\mathbb P_{1,n,v})
<\frac12.
\]

For any normalized finite probability prior \(\pi\), write its complete-record mixture as
\[
\mathbb P_{\pi,n}=\int P^{\otimes n}\,\pi(dP).
\]
The bounded-prior conclusion of the cited result gives a threshold \(N(v)\in\mathbb N\) such that, for every \(n\ge N(v)\) and every pair of normalized finite probability priors \(\pi_0,\pi_1\),
\[
\left.
\begin{aligned}
\pi_0(\{P\})>0&\ \Longrightarrow\ P\in H_0^{\mathrm b}(w),\\
\pi_1(\{P\})>0&\ \Longrightarrow\
P\in\mathcal M_w^{\mathrm b},\quad c_v\rho_n(v)\le d(P)
\end{aligned}
\right\}
\ \Longrightarrow\
\operatorname{TV}(\mathbb P_{\pi_0,n},\mathbb P_{\pi_1,n})
\ge\frac12.
\]
These are precisely the asserted support, magnitude, and complete-record mixture conclusions.
% lean: CausalSmith.Stat.FinitepHomogeneityDensegamma.frontier_at

\item For \(p<2\), the lower ratio bound and the strictly positive exponent difference give
\[
0<\frac{c_v}{C_w^{\mathrm b}},
\qquad
\frac{c_v}{C_w^{\mathrm b}}n^{E(2)-E(p)}
\longrightarrow\infty,
\qquad
\Delta_n(v)\longrightarrow\infty.
\]
The definition in \cref{obj:def:tail-region} includes admissibility and \(p<2\). The preceding divergence supplies its remaining condition, so
\[
\mathcal T=\{v\in\mathcal V:p<2\}.
\]

To characterize \cref{obj:def:equality-region}, suppose that an admissible tuple has a bounded radius ratio. Thus there exists \(C_\Delta\in\mathbb R\) such that
\[
\Delta_n(v)\le C_\Delta
\qquad(n\ge2).
\]
If \(p\ne2\), admissibility implies \(p<2\), and the divergence just established supplies \(N_\Delta\in\mathbb N\) with
\[
n\ge N_\Delta\ \Longrightarrow\ \Delta_n(v)>C_\Delta.
\]
At
\[
n_\Delta=\max\{2,N_\Delta\},
\]
the two bounds would give
\[
C_\Delta<\Delta_{n_\Delta}(v)\le C_\Delta,
\]
a contradiction. Thus boundedness implies \(p=2\). Conversely, when \(p=2\), evaluation at \(2\) preserves the entire tuple, and the upper ratio bound reduces to
\[
\Delta_n(v)\le\frac{C_v}{c_w^{\mathrm b}}
\qquad(n\ge2).
\]
Consequently,
\[
\mathcal E=\{v\in\mathcal V:p=2\}.
\]
% lean: CausalSmith.Stat.FinitepHomogeneityDensegamma.radiusRatio_tendsto_atTop
% lean: CausalSmith.Stat.FinitepHomogeneityDensegamma.tailRegion_characterization
% lean: CausalSmith.Stat.FinitepHomogeneityDensegamma.second_moment_radiusRatio_bounded
% lean: CausalSmith.Stat.FinitepHomogeneityDensegamma.equalityRegion_characterization

\item Define the ambient set
\[
\mathcal U
=\{(p,\alpha,\beta,\gamma)\in\mathbb R^4:p<2\}.
\]
The coordinate map to \(p\) is continuous, so \(\mathcal U\) is open in the product topology. The tuple
\[
v_{\mathrm o}=(3/2,1/2,1/2,1/2)
\]
satisfies every admissibility inequality and belongs to \(\mathcal U\). Thus
\[
v_{\mathrm o}\in\mathcal U\cap\mathcal V,
\qquad
\mathcal U\cap\mathcal V\ne\varnothing.
\]
The characterization of \(\mathcal T\) gives
\[
\mathcal U\cap\mathcal V\subseteq\mathcal T,
\]
which supplies the required open set.
% lean: CausalSmith.Stat.FinitepHomogeneityDensegamma.tailRegion_nonempty_open

\item It remains to prove compact uniformity. For each admissible tuple define the auxiliary positive base and lower envelope by
\[
\Theta(v)
=\max\left\{2,32^{D_p\gamma/(2S)}\right\}+1,
\]
\[
\underline c_{\mathrm{env}}(v)
=\min\left\{
c_0^{\mathrm e},\,
k_v,\,
c_0^{\mathrm e}\Theta(v)^{-|E_0(p)-E_4(p)|}
\right\}.
\]
Every entry in this minimum is positive, so
\[
\underline c_{\mathrm{env}}(v)>0.
\]
The ceiling bound gives
\[
2\le N_{\mathrm{low}}(v)\le\Theta(v).
\]
Since the exponent \(-|E_0(p)-E_4(p)|\) is nonpositive, and since
\(N_{\mathrm{low}}(v)\ge1\), monotonicity in the base and then in the exponent gives
\[
\begin{aligned}
\Theta(v)^{-|E_0(p)-E_4(p)|}
&\le N_{\mathrm{low}}(v)^{-|E_0(p)-E_4(p)|}\\
&\le N_{\mathrm{low}}(v)^{-(E_0(p)-E_4(p))}.
\end{aligned}
\]
On \(F(p)\ge1\), the envelope is at most
\(c_0^{\mathrm e}=c_v\). On \(F(p)<1\), it is at most each of
\[
k_v,
\qquad
c_0^{\mathrm e}
N_{\mathrm{low}}(v)^{-(E_0(p)-E_4(p))}.
\]
The branchwise definition therefore gives, throughout the admissible domain,
\[
0<\underline c_{\mathrm{env}}(v)\le c_v.
\]

The envelope is continuous at every admissible tuple. Indeed, the coordinate maps are continuous; the denominators defining \(q,D_p,E_0,E_4\) are nonzero there; and every variable power in the envelope has a positive base. Absolute value, maximum, minimum, and the remaining arithmetic operations preserve continuity.

Let \(\mathcal K\subseteq\mathcal V\) be compact and nonempty. The continuous envelope attains a minimum at some \(v_{\mathcal K}\in\mathcal K\). Define
\[
c_{\mathcal K}
=\underline c_{\mathrm{env}}(v_{\mathcal K})
=\min_{v\in\mathcal K}\underline c_{\mathrm{env}}(v)>0.
\]
Then
\[
c_{\mathcal K}
\le\underline c_{\mathrm{env}}(v)
\le c_v
\qquad(v\in\mathcal K).
\]
The compact-uniformity conclusion of \cref{obj:thm:original-record-attainable-rate} supplies real constants \(C_{\mathcal K},N_{\mathcal K}\) such that
\[
C_v\le C_{\mathcal K},
\qquad
N_v\le N_{\mathcal K}
\qquad(v\in\mathcal K).
\]
For the empty compact set, the choices
\[
c_{\mathcal K}=1,\qquad C_{\mathcal K}=0,\qquad N_{\mathcal K}=0
\]
satisfy all requirements vacuously.

Finally, let \(\mathcal K_{\mathrm b}\subseteq\mathcal W\) be compact. Its embedded parameter set is
\[
\mathcal K_2
=\{(2,\alpha,\beta,\gamma):
(\alpha,\beta,\gamma)\in\mathcal K_{\mathrm b}\}.
\]
The embedding \(w\mapsto(2,w)\) is continuous, so \(\mathcal K_2\) is compact, and all its tuples are admissible. Applying the preceding result to \(\mathcal K_2\), and using the public identities
\[
c_w^{\mathrm b}=c_{(2,w)},
\qquad
C_w^{\mathrm b}=C_{(2,w)},
\qquad
N_w^{\mathrm b}=N_{(2,w)},
\]
gives
\[
c_{\mathcal K_2}\le c_w^{\mathrm b},
\qquad
C_w^{\mathrm b}\le C_{\mathcal K_2},
\qquad
N_w^{\mathrm b}\le N_{\mathcal K_2}
\qquad(w\in\mathcal K_{\mathrm b}),
\]
with \(c_{\mathcal K_2}>0\). This establishes both compact-uniformity assertions.
% lean: CausalSmith.Stat.FinitepHomogeneityDensegamma.compactLowerEnvelope_le
% lean: CausalSmith.Stat.FinitepHomogeneityDensegamma.continuousAt_compactLowerEnvelope
% lean: CausalSmith.Stat.FinitepHomogeneityDensegamma.cLower_compact_uniformity
% lean: CausalSmith.Stat.FinitepHomogeneityDensegamma.frontier_compact_uniformity
\end{enumerate}
\end{proof}
